\documentclass[12pt]{article}
\usepackage[utf8]{inputenc}
\usepackage[T1]{fontenc}
\usepackage[margin=1in]{geometry}
\usepackage{setspace}
\usepackage{xcolor}
\usepackage[colorlinks=true, urlcolor=blue!60!black, linkcolor=blue!60!black]{hyperref}
\setlength{\parindent}{0em}
\setlength{\parskip}{0.45em}
\renewcommand{\baselinestretch}{1.0}
\begin{document}
\begin{center}
{\Large \textit{Lawrence v. Texas}, 539 U.S. 558 (2003)\par}
\vspace{0.3em}
{\normalsize Case Brief --- One-Page Edition\par}
\end{center}

\textbf{Memory Jogger.} \textit{Bowers} overruled: no criminal ban on private, consensual same-sex intimacy; moral disapproval is not enough (\href{https://www.courtlistener.com/opinion/130160/lawrence-v-texas/}{CourtListener}).

\textbf{Facts.} Deputies answering a false weapons report entered John Lawrence's Houston apartment and found him with Tyron Garner in a consensual sexual act. Both were convicted under Texas's ``Homosexual Conduct'' law, which punished same-sex but not opposite-sex sodomy. Though rarely enforced, a conviction carried a permanent record and sex-offender registration in some states.

\textbf{Procedural History.} Each was fined \$200; the en banc Texas Court of Appeals affirmed under \textit{Bowers v. Hardwick}.

\textbf{Votes.} Reversed 6--3. Kennedy wrote for the Court; O'Connor concurred in the judgment; Scalia (with Rehnquist and Thomas) and Thomas dissented.

\textbf{Holding.} Due process protects consenting adults' private, noncommercial intimate conduct; moral disapproval alone is not a legitimate state interest, and \textit{Bowers} is overruled.

\textbf{Reasoning.} \textit{Bowers} misframed the right as one ``to engage in sodomy''; early sodomy laws did not single out same-sex couples, and history is ``the starting point but not in all cases the ending point.'' The law ``furthers no legitimate state interest which can justify its intrusion into the personal and private life of the individual.'' \textit{Casey} and \textit{Romer} had eroded \textit{Bowers}, and no one had relied on it. O'Connor, applying ``a more searching form of rational basis review,'' found that singling out same-sex conduct violated equal protection. Scalia accused the Court of ``an unheard-of form of rational-basis review'' and warned of ``the end of all morals legislation''; Thomas found no general right of privacy.

\textbf{Examples.} \emph{Same-Side:} a fornication ban falls. \emph{Opposite-Side:} prostitution and sex with minors remain punishable. \emph{Fence-Sitter:} adult incest or polygamy.

\textbf{Critique.} Originalists see ``judicial hubris,'' and \textit{Dobbs} revived \textit{Glucksberg}; Laurence Tribe praised the ``double helix'' of liberty and equality that led to \textit{Obergefell}, while Cass Sunstein noted its unclear standard of review.

\textbf{Key Quotation.} ``They knew times can blind us to certain truths and later generations can see that laws once thought necessary and proper in fact serve only to oppress. As the Constitution endures, persons in every generation can invoke its principles in their own search for greater freedom.'' --- Justice Kennedy, 539 U.S. at 578--79.

\end{document}
